% ============================================================
% CHAPTER 10 — PHYSICAL PLANT MODELING IN SIMULINK
% Parameter Estimation Techniques for Lithium-Ion Batteries
% ============================================================

\chapter{Physical Plant Modeling in Simulink}
\label{ch:simulink_plant}

\section{Introduction to Modeling Approaches}
\label{sec:10_intro}

The preceding chapters treated parameter identification largely as a signal-processing problem: given a record of current, voltage, and temperature, recover the coefficients of an equivalent circuit or an observer gain that best explains the data. This chapter addresses the complementary problem of \emph{producing} that data in a controlled, repeatable, and instrumented setting before any embedded estimator is exposed to a physical cell or to a full vehicle test bench. 

The physical plant model built in this chapter is not the identification algorithm; it is the surrogate battery against which offline and online identification routines from Chapters~4 through~9 are exercised, tuned, and stress-tested. Furthermore, generating high-fidelity synthetic data is a prerequisite for training advanced data-driven estimators, such as Physics-Informed Neural Networks (PINNs) that map thermal and electrical states, prior to hardware deployment. The distinction matters throughout the chapter: the \emph{equivalent circuit model (ECM)} described in Chapter~2 supplies the mathematical structure of the plant, while the \emph{Simulink implementation} described here supplies the numerical environment, current excitation, sensor emulation, and thermal coupling required to generate realistic terminal measurements.

Three considerations motivate carrying out this step in Simulink specifically, rather than in a plain MATLAB script or a lower-level Simscape physical-network description. First, block-diagram simulation exposes the signal-flow structure of the plant explicitly, which is convenient both for pedagogical purposes and for later insertion of an estimator block operating in parallel. Second, Simulink's fixed-step discrete solvers allow the plant to be executed at the same sample rate as an embedded battery management system (BMS). Third, the block library approach scales naturally from a single RC branch (Section~\ref{sec:10_2}) to coupled electro-thermal representations (Section~\ref{sec:10_3_2}) without requiring a change of simulation paradigm.

\subsection{The Plant Model as a Distinct Entity from the Identification Model}

A recurring source of confusion in battery systems engineering is the conflation of the plant with the model used to identify it. In this report, the two are kept deliberately separate:

\begin{itemize}
    \item The \textbf{plant} is the Simulink subsystem intended to stand in for the physical cell. It is parameterized once, from a reference dataset or high-fidelity literature source, and is thereafter treated as ``ground truth.''
    \item The \textbf{identification model} is the structure (for example, an Extended Kalman Filter, an R-GCN, or a PINN) that observes the plant's terminal voltage and current and attempts to recover the plant's parameters, possibly under a different or deliberately mismatched structural assumption.
\end{itemize}

This separation allows a controlled study of structural mismatch, noise sensitivity, and convergence behavior. Committing an ``inverse crime''—using the exact same equations and noise-free conditions to generate and recover data—creates a false sense of algorithmic robustness. 

\subsection{Chapter Roadmap}

Section~\ref{sec:10_1} develops the top-level simulation architecture: the current-source excitation, the SOC integrator, and the noise injection stage. Section~\ref{sec:10_2} works through the parameterization of a first-order Thevenin plant in detail, including the discretization of the RC branch, hysteresis considerations, and lookup table implementation. Section~\ref{sec:10_3} builds the excitation side of the plant using the Dynamic Stress Test (DST) framework. Section~\ref{sec:10_4} summarizes the chapter and connects the plant model to the Software-in-the-Loop validation.

% ============================================================
\section{Simulating the Battery Environment}
\label{sec:10_1}

\subsection{Top-Level Architecture}

The plant model is organized as four coupled subsystems: an excitation subsystem that supplies a current reference $i(t)$; an electrical subsystem that maps $i(t)$ and the present state of charge onto a terminal voltage $V_t(t)$; a state-of-charge integrator via coulomb counting; and a thermal subsystem that accepts the dissipated power and returns a cell temperature. A fifth, non-physical block injects sensor noise and quantization.

This layered architecture ensures every quantity that depends algebraically on SOC (open-circuit voltage, resistances, capacitances) is computed from the integrator output rather than fed back through an implicit loop. This strictly avoids algebraic loops, keeping the model compatible with both variable-step and fixed-step discrete solvers.

\begin{table}[htbp]
\centering
\caption{Comparison of plant-modeling platforms available within the MATLAB/Simulink ecosystem.}
\label{tab:10_platforms}
\renewcommand{\arraystretch}{1.3} % Adds vertical padding
\begin{tabularx}{\textwidth}{@{}l X X@{}}
\toprule
\textbf{Platform} & \textbf{Strengths} & \textbf{Limitations} \\
\midrule
Plain MATLAB script & Full transparency of every computation; trivial to vectorize for Monte Carlo studies; fast batch processing. & No native block-diagram reuse; harder to sync with a discrete BMS control loop; manual solver step sizing. \\
Simulink block diagram & Modular signal flow; native discrete solvers matching embedded sample rates; easy SIL testing. & Higher per-step overhead; care required to avoid unintended algebraic loops in nonlinear tables. \\
Simscape Electrical & Automatically enforces conservation laws; excellent for multi-cell packs and thermal networks. & Acausal solver can be slower for single-cell studies; physical-network formulation is less transparent for state-space teaching. \\
\bottomrule
\end{tabularx}
\end{table}

\subsection{Solver Selection and Configuration}

Two practical choices strongly affect whether a Simulink plant is fit for identification studies: solver type and step size. A fixed-step discrete solver (such as \texttt{ode3} or Simulink's \texttt{discrete} solver) operating at the sample interval $\Delta t$ is preferred over a variable-step continuous solver (like \texttt{ode45}). It avoids solver-dependent artifacts where an adaptive step size silently refines its resolution around a current step edge, producing transient behavior invisible to a real digital controller.

The step size must be chosen relative to the fastest time constant retained in the electrical subsystem. For a first-order Thevenin branch with polarization resistance $R_1$ and capacitance $C_1$, the branch time constant is $\tau_1 = R_1 C_1$, leading to the heuristic:
\begin{equation}
\Delta t \le \frac{\tau_{1,\min}}{5}
\label{eq:10_step_rule}
\end{equation}
where $\tau_{1,\min}$ is the smallest branch time constant observed across the SOC and temperature operating range. 

% ============================================================
\section{First-Order Thevenin Parameterization}
\label{sec:10_2}

\subsection{Governing Equations}

The plant developed in this section follows the first-order Thevenin structure. Under the sign convention in which discharge current is taken as positive, the continuous-time state equation for the single RC (polarization) branch is:
\begin{equation}
\frac{dU_1}{dt} = -\frac{U_1}{R_1 C_1} + \frac{i(t)}{C_1}
\label{eq:10_rc_state}
\end{equation}
and the terminal voltage is obtained from Kirchhoff's voltage law as:
\begin{equation}
V_t(t) = V_{oc}\big(z(t), T(t)\big) - U_1(t) - i(t) R_0\big(z(t), T(t)\big)
\label{eq:10_terminal}
\end{equation}
where $U_1$ is the voltage across the polarization branch, $R_0$ is the ohmic resistance, $R_1$ and $C_1$ are the polarization resistance and capacitance, $V_{oc}$ is the open-circuit voltage, $z(t)$ is the state of charge, and $T(t)$ is the cell temperature. 

\begin{figure}[htbp]
\centering
% TIKZ CODE FOR CIRCUIT
\begin{circuitikz}[scale=1.2, transform shape]
    % Voltage source (Voc)
    \draw (0,0) to[battery1, l=$V_{oc}(z{,}T)$] (0,3);
    % Ohmic resistance
    \draw (0,3) to[R, l=$R_0(z{,}T)$, i=$i(t)$] (3,3);
    % RC Branch
    \draw (3,3) -- (4,3);
    \draw (4,3) to[R, l=$R_1$] (4,1.5) -- (4,0);
    \draw (3,3) -- (5,3);
    \draw (5,3) to[C, l=$C_1$] (5,1.5) -- (5,0);
    \draw (4,0) -- (3,0);
    \draw (5,0) -- (4,0);
    % Terminals
    \draw (5,3) -- (6.5,3) node[ocirc, label=above:$+$]{};
    \draw (0,0) -- (6.5,0) node[ocirc, label=below:$-$]{};
    % Terminal Voltage Label
    \draw[<->, >=stealth] (6.5,2.8) -- (6.5,0.2) node[midway, fill=white] {$V_t(t)$};
    % Branch Voltage Label
    \draw[<->, >=stealth] (3.6,2.8) -- (3.6,0.2) node[midway, fill=white] {$U_1(t)$};
\end{circuitikz}
\caption{First-order Thevenin equivalent circuit model forming the basis of the electrical subsystem.}
\label{fig:10_thevenin}
\end{figure}

State of charge evolves by coulomb counting:
\begin{equation}
\frac{dz}{dt} = -\frac{\eta \, i(t)}{Q_n}
\label{eq:10_soc}
\end{equation}
with $Q_n$ the nominal capacity and $\eta$ a coulombic efficiency factor.

\subsection{Discrete-Time Implementation}

Because the Simulink plant is executed with a fixed-step discrete solver, equation \eqref{eq:10_rc_state} is discretized exactly, treating the input current as piecewise constant over each sample interval (a zero-order-hold assumption):
\begin{equation}
U_{1,k+1} = e^{-\Delta t / \tau_1} \, U_{1,k} + R_1\left(1 - e^{-\Delta t/\tau_1}\right) i_k, \qquad \tau_1 = R_1 C_1
\label{eq:10_rc_discrete}
\end{equation}
The state-of-charge integrator is discretized identically:
\begin{equation}
z_{k+1} = z_k - \frac{\eta \, \Delta t}{Q_n} \, i_k .
\label{eq:10_soc_discrete}
\end{equation}

\subsection{Embedding SOC-Dependent Parameters}

$V_{oc}$, $R_0$, $R_1$, and $C_1$ are all functions of SOC and temperature. Two implementation strategies are in common use and are compared in Table~\ref{tab:10_lookup_vs_poly}.

\begin{table}[htbp]
\centering
\caption{Table-based versus polynomial implementation of SOC-dependent ECM parameters in Simulink.}
\label{tab:10_lookup_vs_poly}
\renewcommand{\arraystretch}{1.3}
\begin{tabularx}{\textwidth}{@{}l X X@{}}
\toprule
\textbf{Criterion} & \textbf{Lookup table (1-D/2-D)} & \textbf{Polynomial curve-fit} \\
\midrule
Fidelity near extremes & Preserves sharp measured transitions if grid is dense enough. & Can smooth over genuine nonlinearities unless high order is used. \\
Extrapolation & Clamped outside grid range; must be guarded explicitly. & Extrapolates automatically, but with potentially non-physical values. \\
Run-time cost & One interpolation lookup per parameter per step. & One polynomial evaluation per parameter per step. \\
Traceability & Directly traceable to individual HPPC measurements. & Coefficients are derived; individual points are not directly recoverable. \\
\bottomrule
\end{tabularx}
\end{table}

\begin{table}[htbp]
\centering
\caption{Illustrative parameters for an NMC~811/graphite cell at 1C discharge rate. (Source: Barletta et al. \cite{barletta2022})}
\label{tab:10_barletta}
\begin{tabular}{ccccc}
\toprule
SOC (\%) & $V_{oc}$ (V) & $R_0$ (m$\Omega$) & $R_1$ (m$\Omega$) & $C_1$ (kF) \\
\midrule
20  & 3.853 & 9.143 & 0.472 & 58.55 \\
60  & 3.979 & 6.112 & 0.399 & 68.25 \\
100 & 4.157 & 2.315 & 0.201 & 290.9 \\
\bottomrule
\end{tabular}
\end{table}

% ============================================================
\section{Dynamic Stress Testing Framework}
\label{sec:10_3}

\subsection{Purpose and Provenance of the DST Profile}

While HPPC tests characterize parameters at quasi-static SOC points, the Dynamic Stress Test (DST) exercises the parameterized model under conditions resembling a vehicle drive cycle. The DST originates from the USABC/PNGV manuals, derived from the Simplified Federal Urban Driving Schedule (SFUDS). 

\begin{table}[htbp]
\centering
\small
\caption{Distinguishing features of the principal excitation profiles.}
\label{tab:10_profiles}
\renewcommand{\arraystretch}{1.3}
\begin{tabularx}{\textwidth}{@{}l >{\raggedright\arraybackslash}X >{\raggedright\arraybackslash}X l@{}}
\toprule
\textbf{Profile} & \textbf{Primary Use} & \textbf{Excitation Character} & \textbf{Source} \\
\midrule
HPPC & Quasi-static char. & Discrete pulses at rest. & USABC/FreedomCAR \\
DST & Dynamic ID \& validation & Variable-power micro-cycle. & USABC/PNGV \\
FUDS & Out-of-sample validation & Recorded drive trace. & Federal standards \\
\bottomrule
\end{tabularx}
\end{table}

\subsection{High-Current Stress and Thermal Boundary Mapping}
\label{sec:10_3_2}

The high-magnitude pulses in a DST profile produce non-negligible self-heating. The heat generation rate is computed using the Bernardi energy-balance form, decomposing total heat into an irreversible (Joule) term and a reversible (entropic) term:
\begin{equation}
\dot Q(t) = i(t)^2 \Big[R_0(z,T) + R_1(z,T)\Big] \;-\; i(t)\, T(t) \, \frac{\partial V_{oc}}{\partial T}
\label{eq:10_bernardi}
\end{equation}

\begin{algorithm}[htbp]
\caption{Simulated thermal boundary mapping using the coupled plant}
\label{alg:10_boundary_map}
\begin{algorithmic}[1]
\State \textbf{Inputs:} Candidate current magnitudes $\{i_1,\dots,i_m\}$; candidate pulse durations $\{d_1,\dots,d_n\}$; ambient temperatures $\{T_{a,1},\dots,T_{a,p}\}$; max temperature rise $\Delta T_{lim}$.
\State \textbf{Outputs:} Map of simulated temp rise $\Delta T(i, d, T_a)$; boundary current $i_{lim}(d, T_a)$.
\For{each ambient temperature $T_{a,k}$}
    \State Initialize plant at a fixed reference SOC and set $T(0) = T_{a,k}$
    \For{each pulse duration $d_j$}
        \For{each current magnitude $i_\ell$}
            \State Apply discharge pulse of magnitude $i_\ell$ for duration $d_j$.
            \State Log $\Delta T = T(d_j) - T_{a,k}$ from the thermal subsystem.
            \State Record $\dot Q(t)$ from equation \eqref{eq:10_bernardi} for diagnostics.
        \EndFor
        \State Interpolate across $\{i_\ell\}$ to estimate $i_{lim}(d_j, T_{a,k})$ such that $\Delta T = \Delta T_{lim}$.
    \EndFor
\EndFor
\State \Return tabulated map $\Delta T$ and boundary curve $i_{lim}$
\end{algorithmic}
\end{algorithm}

% ============================================================
\section{Summary}
\label{sec:10_4}

This chapter converted the equivalent circuit structures and offline characterization procedures into an executable Simulink plant. The plant was organized around a clean separation between excitation, electrical, state-of-charge, and thermal subsystems. The use of strict zero-order-hold discretization and deliberate algebraic loop avoidance ensures the plant runs robustly alongside complex digital observers, paving the way for testing advanced architectures like adaptive filters or Physics-Informed Neural Networks (PINNs) in subsequent Software-in-the-Loop implementations.

% ============================================================
% AI IMAGE GENERATION PROMPTS
% ============================================================
% If you wish to include realistic visualizations, use the following 
% prompts in an AI image generator (like Midjourney or DALL-E 3):
%
% PROMPT 1 (For Section 10.1 - Hardware representation):
% "A photorealistic macro shot of a cylindrical lithium-ion battery 
% connected to a high-end laboratory testing rig. Thick copper wires 
% with precise sensory clamps are attached to the battery terminals. 
% The background features out-of-focus digital screens displaying 
% real-time data charts and Simulink block diagrams. Studio lighting, 
% high technological aesthetic."
%
% PROMPT 2 (For Section 10.3 - Thermal Stress):
% "A scientific thermal imaging (FLIR style) photograph of a 
% lithium-ion battery pack undergoing a high-current stress test. 
% The battery cells glow in gradients of bright yellow and red, 
% indicating heat generation, against a dark blue background. High 
% contrast, technical and scientific visual style."
% ============================================================